\documentclass[10pt,a4paper]{amsart}
\usepackage[margin=19mm]{geometry}
\usepackage{amsmath,amssymb,mathtools}
\usepackage[scr=rsfs,cal=euler]{mathalfa}
\usepackage{xfrac}
\usepackage[unicode,hidelinks,bookmarks=false]{hyperref}
\newcommand{\ie}{i.e.\ }
\newcommand{\cf}{cf.\ }
\newcommand{\resp}{respectively\ }
\newcommand{\Subsection}{\subsection}
\providecommand{\nobreakdash}{\nobreak\hbox{-}}
\setlength{\emergencystretch}{2em}
\multlinegap0pt
\usepackage{bm}
\usepackage{bbm}
\usepackage{dsfont}
\usepackage{xspace}
\usepackage{cancel}
\usepackage{mathtools}
\usepackage{amsthm}
\makeatletter
\@ifundefined{lemma}{\newtheorem{lemma}{Lemma}}{}
\@ifundefined{proposition}{\newtheorem{proposition}{Proposition}}{}
\makeatother
\DeclareMathOperator{\BetaDist}{Beta}
\title[Random-Subspace Frank--Wolfe over Strongly Convex Sets]{Random-Subspace Frank--Wolfe over Strongly Convex Sets}
\author{Pierre-Louis Poirion, Sebastian Pokutta, Akiko Takeda}
\date{Source: arXiv:2605.24819v1 (CC BY 4.0)}
\begin{document}
\maketitle

\noindent\emph{Source and licence.} Random-Subspace Frank--Wolfe over Strongly Convex Sets. Version arXiv:2605.24819v1 (CC BY 4.0), distributed under the Creative Commons Attribution 4.0 International licence, as recorded in the arXiv licence metadata for this exact version and shown on the abstract page https://arxiv.org/abs/2605.24819v1. Full rights notice: licenses/arxiv-2605.24819-CC-BY-4.0.txt.\par\smallskip

\noindent\textit{Editorial scope.} Counted unit: the Proposition whose source label is \texttt{prop:exact\_Delta2\_over\_Y} in the article (source0.tex lines 2087--2099, source label \texttt{prop:exact\_Delta2\_over\_Y}), delivered together with the article's own complete original proof of it (source0.tex lines 2101--2369, one \texttt{proof} environment of 269 lines). No mathematics is written, repaired or re-derived: statement and proof are the article's own text line for line. Required context retained uncounted: lem:cond\_decomp (lines 1978--2023); eq:T\_moments (lines 1999--2006). Every definition, display and label of those lines is retained unchanged. Omitted from this package and not redistributed: 1--1977, 2007--2086, 2100--2100, 2370--6207. Nothing inside a retained excerpt is omitted or altered: every retained body is delivered exactly as the article's lines are written, so no omission receipt of any kind accompanies this package. Citations: the retained excerpts carry no citation command at all, so no bibliography entry of the article is transcribed and no reference is redistributed. Reference adaptation: none was needed and none was made; every reference inside the delivered unit resolves inside it, so no unresolved reference or citation is produced. Printed slips kept literal: no printed word, symbol, hypothesis or number of the counted statement or of its proof was altered. Layout and class: the article's own class is \texttt{article} with packages including neurips\_2026, amsmath, amsthm, amssymb, mathtools, algorithm, algpseudocode, microtype, hyperref, enumitem, graphicx, xcolor, caption, subcaption; the wrapper is the lane's pinned \texttt{amsart} editorial wrapper at 10pt on A4, which restates the theorem-like environments the retained text needs and adds the article's own macro definitions that the retained text needs (\texttt{\textbackslash BetaDist}), with extra wrapper packages (bm, bbm, dsfont, xspace, cancel, mathtools). The delivered numbering is the unit's own. Assets: no retained span contains a figure environment, a graphics-inclusion command, a PDF-page inclusion, a table or a TikZ picture; no image file is copied into this package, the package declares no asset dependency and no image file is redistributed with it. Licence: version arXiv:2605.24819v1 of this article is distributed under the Creative Commons Attribution 4.0 International licence, as recorded in the arXiv licence metadata for this exact version and shown on the abstract page https://arxiv.org/abs/2605.24819v1. A full rights notice is delivered at licenses/arxiv-2605.24819-CC-BY-4.0.txt. Characters: every retained excerpt is pure ASCII, so the delivered engine, XeLaTeX with its default Latin Modern text font, reports no missing character.
\begin{lemma}[Exact conditional decomposition]
\label{lem:cond_decomp}
By rotational invariance, assume without loss of generality that
\begin{equation}
v=e_1,
\qquad
u=\rho e_1+\sqrt{1-\rho^2}\,e_2.
\label{eq:wlog_basis}
\end{equation}
Set
\[
\sigma:=\sqrt{1-\rho^2}.
\]
Assume \(2\le d<n\), then there exist random variables $(S,T,B)$ such that:
\begin{enumerate}
\item $S\in(0,1)$ almost surely and
\begin{equation}
S\sim \BetaDist\!\left(\tfrac d2,\tfrac{n-d}{2}\right).
\label{eq:S_beta}
\end{equation}
\item $T\in[-1,1]$ is independent of $S$ and satisfies
\begin{equation}
\mathbb E[T]=0,
\qquad
\mathbb E[T^2]=\frac{1}{n-1},
\qquad
\mathbb E[T^4]=\frac{3}{(n-1)(n+1)}.
\label{eq:T_moments}
\end{equation}
\item \(B\) is independent of \((S,T)\). If \(d\le n-2\), then
\[
B\sim \BetaDist\!\left(\tfrac{d-1}{2},\tfrac{n-d-1}{2}\right),
\]
while if \(d=n-1\), we use the degenerate convention \(B\equiv1\).
\item Defining
\begin{equation}
\alpha:=\rho\sqrt S+\sigma\sqrt{1-S}\,T,
\label{eq:def_alpha}
\end{equation}
one has
\begin{equation}
X=\alpha^2+\sigma^2(1-T^2)B.
\label{eq:X_decomp}
\end{equation}
\end{enumerate}
\end{lemma}

\begin{proposition}[Exact formula for \(\mathbb E((X-Y)^2/Y)\)]
\label{prop:exact_Delta2_over_Y}
Assume \(3\le d<n\). Then
\[
\mathbb E\!\left[\frac{(X-Y)^2}{Y}\right]
=
\frac{4(n-d)}{n(n-1)}(1-\rho^2)
\left(
\rho^2+
\frac{(1-\rho^2)(n-d+2)}{2(d-2)(n+1)}
\right).
\]
\end{proposition}

\begin{proof}
Let us give first a global overview of the proof.
\paragraph{Proof idea:}
The identity for \(\mathbb E[\Gamma]\) reduces the problem to controlling
\(\mathbb E(\|Pu\|-\|Pv\|)^2\). We write
\(X=\|Pu\|^2\), \(Y=\|Pv\|^2\), use
\((\sqrt X-\sqrt Y)^2\le (X-Y)^2/Y\), and then compute the latter exactly using
the conditional decomposition from Lemma~\ref{lem:cond_decomp}. The calculation
is only algebraic: after conditioning on \(S=Y\), all odd powers of \(T\)
vanish and the remaining beta moments are explicit.



Set
\[
\sigma:=\sqrt{1-\rho^2}.
\]
By Lemma~\ref{lem:cond_decomp},
\begin{equation}
X-S=(\alpha^2-S)+\sigma^2(1-T^2)B.
\label{eq:Delta_split_corrected}
\end{equation}
Therefore
\begin{equation}
(X-S)^2
=
(\alpha^2-S)^2
+
2(\alpha^2-S)\sigma^2(1-T^2)B
+
\sigma^4(1-T^2)^2B^2.
\label{eq:expand_square_corrected}
\end{equation}

Conditional on $(S,T)$, the Beta variable
\[
B\sim \BetaDist\!\left(\tfrac{d-1}{2},\tfrac{n-d-1}{2}\right)
\]
is independent of $(S,T)$, hence
\begin{equation}
\mathbb E\big[(X-S)^2\mid S,T\big]
=
(\alpha^2-S)^2
+
2(\alpha^2-S)\sigma^2(1-T^2)\,\mathbb E[B]
+
\sigma^4(1-T^2)^2\,\mathbb E[B^2].
\label{eq:cond_second_moment_corrected}
\end{equation}
For this Beta distribution,
\begin{equation}
\mathbb E[B]=\frac{d-1}{n-2},
\qquad
\mathbb E[B^2]=\frac{d^2-1}{n(n-2)}.
\label{eq:B_moments_corrected}
\end{equation}

We now average over $T$ conditional on $S$.
Using
\[
\alpha=\rho\sqrt S+\sigma\sqrt{1-S}\,T,
\]
we get
\begin{equation}
\alpha^2-S
=
\sigma^2\bigl(-S+(1-S)T^2\bigr)
+
2\rho\sigma\sqrt{S(1-S)}\,T.
\label{eq:alpha2_minus_S_corrected}
\end{equation}
Define
\[
A(S,T):=\sigma^2\bigl(-S+(1-S)T^2\bigr),
\qquad
L(S):=2\rho\sigma\sqrt{S(1-S)},
\]
so that
\[
\alpha^2-S=A(S,T)+L(S)\,T.
\]

Since $T$ is symmetric and $\mathbb E[T]=\mathbb E[T^3]=0$, all odd terms in $T$ vanish after taking expectation.
Using \eqref{eq:T_moments}, we obtain
\begin{align}
\mathbb E\big[(\alpha^2-S)^2\mid S\big]
&=
\mathbb E\big[A(S,T)^2\mid S\big]
+
\mathbb E\big[L(S)^2T^2\mid S\big]
\notag\\
&=
\sigma^4\,
\mathbb E\!\left[\bigl(-S+(1-S)T^2\bigr)^2\mid S\right]
+
4\rho^2\sigma^2 S(1-S)\,\mathbb E[T^2]
\notag\\
&=
\sigma^4\left(
S^2
-2S(1-S)\mathbb E[T^2]
+(1-S)^2\mathbb E[T^4]
\right)
+
\frac{4\rho^2\sigma^2}{n-1}S(1-S)
\notag\\
&=
\sigma^4\left(
S^2
-\frac{2S(1-S)}{n-1}
+\frac{3(1-S)^2}{(n-1)(n+1)}
\right)
+
\frac{4\rho^2\sigma^2}{n-1}S(1-S).
\label{eq:ET_alphaS_sq_corrected}
\end{align}

Similarly,
\begin{align}
\mathbb E\big[(\alpha^2-S)(1-T^2)\mid S\big]
&=
\mathbb E\big[A(S,T)(1-T^2)\mid S\big]
\notag\\
&=
\sigma^2\,\mathbb E\!\left[\bigl(-S+(1-S)T^2\bigr)(1-T^2)\mid S\right]
\notag\\
&=
\sigma^2\left(
-S(1-\mathbb E[T^2])+(1-S)(\mathbb E[T^2]-\mathbb E[T^4])
\right)
\notag\\
&=
\sigma^2\left(
-S\frac{n-2}{n-1}
+
(1-S)\frac{n-2}{(n-1)(n+1)}
\right).
\label{eq:ET_cross_corrected}
\end{align}

Finally,
\begin{align}
\mathbb E\big[(1-T^2)^2\big]
&=
1-2\mathbb E[T^2]+\mathbb E[T^4]
\notag\\
&=
1-\frac{2}{n-1}+\frac{3}{(n-1)(n+1)}
\notag\\
&=
\frac{n(n-2)}{(n-1)(n+1)}.
\label{eq:ET_one_minus_T2_sq_corrected}
\end{align}

Substituting \eqref{eq:B_moments_corrected}, \eqref{eq:ET_alphaS_sq_corrected}, \eqref{eq:ET_cross_corrected}, and \eqref{eq:ET_one_minus_T2_sq_corrected} into \eqref{eq:cond_second_moment_corrected}, we get
\begin{align}
\mathbb E\big[(X-S)^2\mid S\big]
&=
\sigma^4\left(
S^2
-\frac{2S(1-S)}{n-1}
+\frac{3(1-S)^2}{(n-1)(n+1)}
\right)
+
\frac{4\rho^2\sigma^2}{n-1}S(1-S)
\notag\\
&\quad
+
2\sigma^2
\left[
\sigma^2\left(
-S\frac{n-2}{n-1}
+
(1-S)\frac{n-2}{(n-1)(n+1)}
\right)
\right]
\frac{d-1}{n-2}
\notag\\
&\quad
+
\sigma^4\frac{n(n-2)}{(n-1)(n+1)}\frac{d^2-1}{n(n-2)}.
\label{eq:after_substitution_corrected}
\end{align}

Collecting the $\sigma^4$-terms and simplifying yields
\begin{equation}
\mathbb E\big[(X-S)^2\mid S\big]
=
\frac{4\rho^2\sigma^2}{n-1}S(1-S)
+
\frac{\sigma^4}{(n-1)(n+1)}
\Bigl((n^2+2n+4)S^2-2(dn+2d+2)S+d(d+2)\Bigr).
\label{eq:cond_second_moment_compact_corrected}
\end{equation}

We now divide by $S$ and take expectation over $S$:
\begin{align}
\mathbb E\!\left[\frac{(X-S)^2}{S}\right]
&=
\frac{4\rho^2\sigma^2}{n-1}\,\mathbb E[1-S]
\notag\\
&\quad
+
\frac{\sigma^4}{(n-1)(n+1)}
\left(
(n^2+2n+4)\mathbb E[S]
-2(dn+2d+2)
+d(d+2)\mathbb E\!\left[\frac1S\right]
\right).
\label{eq:Delta2_over_S_split_corrected}
\end{align}

Since
\[
S\sim \BetaDist\!\left(\tfrac d2,\tfrac{n-d}{2}\right),
\]
we have
\begin{equation}
\mathbb E[S]=\frac{d}{n},
\qquad
\mathbb E[1-S]=\frac{n-d}{n},
\qquad
\mathbb E\!\left[\frac1S\right]=\frac{n-2}{d-2},
\quad d>2.
\label{eq:S_moments_corrected}
\end{equation}
Substituting \eqref{eq:S_moments_corrected} into \eqref{eq:Delta2_over_S_split_corrected} gives
\begin{align}
\mathbb E\!\left[\frac{(X-S)^2}{S}\right]
&=
\frac{4\rho^2\sigma^2(n-d)}{n(n-1)}
\notag\\
&\quad
+
\frac{\sigma^4}{(n-1)(n+1)}
\left(
(n^2+2n+4)\frac dn
-2(dn+2d+2)
+d(d+2)\frac{n-2}{d-2}
\right).
\label{eq:after_S_moments_corrected}
\end{align}

The bracket simplifies to
\begin{equation}
(n^2+2n+4)\frac dn
-2(dn+2d+2)
+d(d+2)\frac{n-2}{d-2}
=
\frac{2(n-d)(n-d+2)(n+1)}{n(d-2)}.
\label{eq:bracket_common_den_corrected}
\end{equation}
Substituting \eqref{eq:bracket_common_den_corrected} into \eqref{eq:after_S_moments_corrected}, we obtain
\begin{align}
\mathbb E\!\left[\frac{(X-S)^2}{S}\right]
&=
\frac{4\rho^2\sigma^2(n-d)}{n(n-1)}
+
\frac{2\sigma^4(n-d)(n-d+2)}{n(n-1)(d-2)(n+1)}
\notag\\
&=
\frac{4(n-d)}{n(n-1)}(1-\rho^2)
\left(
\rho^2+\frac{(1-\rho^2)(n-d+2)}{2(d-2)(n+1)}
\right).
\label{eq:almost_done_corrected}
\end{align}
This is the claimed formula.
\end{proof}

\end{document}
